% Authored: reproducibility and data provenance.

\chapter{Reproducibility}
\label{ch:reproducibility}

\section{The reproducibility equation}

A result produced by this repository is determined by six inputs:

\begin{equation}
\underbrace{\text{Environment}}_{\text{Python + poetry.lock}}
\;+\;
\underbrace{\text{Revision}}_{\text{git SHA}}
\;+\;
\underbrace{\text{Configuration}}_{\text{Settings / .env}}
\;+\;
\underbrace{\text{Data}}_{\text{versioned CSV}}
\;+\;
\underbrace{\text{Seeds}}_{\text{local Random}}
\;+\;
\underbrace{\text{Pipeline}}_{\text{CLI or make target}}
\;=\;
\text{Artifact}
\label{eq:reproducibility}
\end{equation}

Each term corresponds to something concrete in the repository:

\begin{description}
  \item[Environment] The supported interpreter range is declared in
    \srcfile{pyproject.toml} and exercised across 3.10, 3.11, and 3.12 in CI.
    Exact package versions are pinned in \srcfile{poetry.lock}, which is
    version controlled. The container image pins Poetry itself.
  \item[Revision] Every artifact is attributable to a git SHA. This manual
    records the SHA it describes on its title page and in the documentation
    manifest.
  \item[Configuration] A single settings object with declared defaults, so the
    effective configuration of a run is the defaults plus the environment
    variables listed in \cref{ch:configuration}.
  \item[Data] The sample dataset and every fixture are authored files under
    version control, so ``the data'' is fully identified by the revision.
  \item[Seeds] Synthetic generation uses a locally seeded
    \texttt{random.Random}; nothing in the package touches the global
    \texttt{random} state or a NumPy generator. Nothing else in the default
    path is stochastic at all.
  \item[Pipeline] Executed through a documented CLI command or Make target.
\end{description}

\section{What is deterministic}

With the default configuration the entire pipeline is deterministic:

\begin{itemize}[nosep]
  \item embedding (hash-derived, no learned parameters);
  \item retrieval and ranking (fixed scoring, ties broken by corpus position);
  \item reranking (fixed linear combination of transparent signals);
  \item generation (the deterministic local provider);
  \item output parsing and repair;
  \item citation construction;
  \item evidence-overlap checking;
  \item evaluation metrics and their aggregation;
  \item synthetic dataset generation for a fixed seed and parameters;
  \item lakehouse export, which the module describes as writing deterministic
    partitioned files;
  \item documentation generation, which is byte-identical for a given revision.
\end{itemize}

\section{What is not, and why}

Stating the boundary precisely matters more than claiming reproducibility:

\begin{description}
  \item[Remote LLM providers] A remote model is not reproducible across time.
    The provider may be updated, and the repository sets no seed for remote
    calls (the Bedrock adapter sets a temperature, which reduces but does not
    remove variability). Any evaluation or experiment run against a remote
    provider is a point-in-time measurement.
  \item[Wall-clock timing] Benchmark latencies depend on the machine. The
    harness reports sample counts and five statistics rather than a single
    number, but the values are not comparable across hosts.
  \item[Timestamps in artifacts] Saved reports, jobs, conversations, and
    telemetry events record creation times, so those files differ between runs
    by construction.
  \item[Dependency resolution outside the lock] Installing without
    \srcfile{poetry.lock} --- or installing extras, which the lock also covers
    but CI never installs --- can resolve differently.
  \item[Retrieval over a changing corpus] Incremental index updates change the
    corpus, and BM25 scores depend on corpus statistics
    ($N$, $n(t)$, $\overline{|d|}$). Adding documents therefore changes lexical
    scores for existing documents. This is a property of BM25, not a defect,
    but it means a retrieval result is reproducible only against a fixed index.
\end{description}

\section{Checksums and manifests}

The repository records no checksum manifest for its datasets or artifacts. The
only digest mechanism present is the one introduced by the documentation
system: \srcfile{docs/metadata/documentation-manifest.json} stores a SHA-256
for every generated documentation file together with the revision it was
generated from, so documentation drift is detectable. Extending the same idea to
data artifacts is listed in \cref{ch:improvements}.

\section{Data provenance}

% !TEX root = ../main.tex
% ---------------------------------------------------------------
% GENERATED FILE - DO NOT EDIT.
% Produced by docs/tools/render_engineering.py from the repository source.
% Regenerate with: make docs
% ---------------------------------------------------------------

\section{Version-controlled datasets}

Every dataset used by demos, tests, and CI is authored inside the repository. The repository records no external data source, no acquisition script, no download URL, and no checksum manifest, so there is no external provenance chain to document.

\begin{longtable}{>{\raggedright\arraybackslash\small}p{0.26\textwidth} >{\raggedright\arraybackslash\small}p{0.08\textwidth} >{\raggedright\arraybackslash\small}p{0.09\textwidth} >{\raggedright\arraybackslash\small}p{0.39\textwidth} >{\raggedright\arraybackslash\small}p{0.10\textwidth}}
\caption{Datasets and example inputs under version control.}\label{tab:datasets}\\
\toprule
Artifact & Format & Records & Fields & Bytes \\
\midrule
\endfirsthead
\toprule
Artifact & Format & Records & Fields & Bytes \\
\midrule
\endhead
\midrule
\multicolumn{5}{r}{\footnotesize continued on next page}\\
\endfoot
\bottomrule
\endlastfoot
\texttt{data/\allowbreak{}sample\_\allowbreak{}feedback.\allowbreak{}csv} & csv & 12 & \texttt{feedback\_\allowbreak{}id,\allowbreak{} customer\_\allowbreak{}segment,\allowbreak{} channel,\allowbreak{} rating,\allowbreak{} text,\allowbreak{} created\_\allowbreak{}at} & 2292 \\
\texttt{examples/\allowbreak{}experiment\_\allowbreak{}config.\allowbreak{}yaml} & yaml & -- & -- & 717 \\
\texttt{examples/\allowbreak{}queries.\allowbreak{}jsonl} & jsonl & 5 & \texttt{expected\_\allowbreak{}keywords,\allowbreak{} is\_\allowbreak{}answerable,\allowbreak{} question,\allowbreak{} relevant\_\allowbreak{}document\_\allowbreak{}ids} & 906 \\
\texttt{examples/\allowbreak{}stream\_\allowbreak{}feedback.\allowbreak{}jsonl} & jsonl & 2 & \texttt{channel,\allowbreak{} created\_\allowbreak{}at,\allowbreak{} customer\_\allowbreak{}segment,\allowbreak{} feedback\_\allowbreak{}id,\allowbreak{} rating,\allowbreak{} text} & 428 \\
\texttt{examples/\allowbreak{}tool\_\allowbreak{}queries.\allowbreak{}jsonl} & jsonl & 9 & \texttt{expected\_\allowbreak{}tool,\allowbreak{} note,\allowbreak{} question} & 1612 \\
\end{longtable}

\section{Artifact lineage}

\begin{longtable}{>{\raggedright\arraybackslash\small}p{0.17\textwidth} >{\raggedright\arraybackslash\small}p{0.18\textwidth} >{\raggedright\arraybackslash\small}p{0.20\textwidth} >{\raggedright\arraybackslash\small}p{0.22\textwidth} >{\raggedright\arraybackslash\small}p{0.17\textwidth}}
\caption{Transformation chain from raw CSV to published artifacts.}\label{tab:lineage}\\
\toprule
Source artifact & Producer & Transformation & Validation performed & Destination artifact \\
\midrule
\endfirsthead
\toprule
Source artifact & Producer & Transformation & Validation performed & Destination artifact \\
\midrule
\endhead
\midrule
\multicolumn{5}{r}{\footnotesize continued on next page}\\
\endfoot
\bottomrule
\endlastfoot
data/sample\_feedback.csv & Authored in the repository (no external source) & \texttt{load\_feedback\_csv} & \texttt{validate\_feedback\_csv} (data contract) & Validated \texttt{FeedbackRecord} list \\
Validated records & \texttt{load\_feedback\_csv} & \texttt{feedback\_to\_chunks} & PII redaction via \texttt{redact\_pii} & \texttt{DocumentChunk} list \\
Document chunks & \texttt{feedback\_to\_chunks} & \texttt{HashingEmbeddingModel.embed} & Dimension check in \texttt{InMemoryVectorStore.add} & .artifacts/vector\_store.json \\
Vector index & \texttt{build\_index} & \texttt{QueryEngine} / \texttt{BM25Retriever} / \texttt{HybridRetriever} & Guardrail context check & Ranked \texttt{SearchResult} list \\
Search results & Retriever & \texttt{FeedbackInsightAgent.answer} & Evidence-overlap hallucination check & \texttt{AgentAnswer} with citations \\
Agent answers & \texttt{FeedbackInsightAgent} & \texttt{evaluate\_system} & Metric aggregation over evaluation cases & \texttt{EvaluationReport} JSON \\
Validated records & \texttt{load\_feedback\_csv} & \texttt{export\_feedback\_lakehouse} & Partition column validation, record redaction & Partitioned JSONL + table metadata \\
Stream events & \texttt{JsonlFeedbackStream} / Kafka / Kinesis & \texttt{consume\_feedback\_stream} & Per-event contract validation & \texttt{StreamIngestionResult} + CSV \\
\end{longtable}

